\section{RIPPLE: Corpus-Anchored Geometry Distillation}
\label{sec:method}

RIPPLE implements the requirement derived in Section~\ref{sec:theory} in four stages: corpus anchoring, multi-scale diffusion, candidate sampling, and cross-anchor row supervision. For a distribution $p$ and a set $C$ with $p(C)>0$, $p|_C$ denotes restriction to $C$ followed by renormalization.

\subsection{Corpus-Level Teacher Graph}
\label{sec:graph}

We cache the teacher embeddings and construct a mutual $k$-nearest-neighbor graph over the corpus. Its support is
\begin{equation}
\label{eq:mutual-support}
j\in\mathcal M_k^T(i)
\quad\Longleftrightarrow\quad
j\in\mathcal N_k^T(i)\ \text{and}\ i\in\mathcal N_k^T(j),
\end{equation}
which retains reciprocal teacher neighborhoods and removes one-directional hub edges. The weighted transition operator is
\begin{equation}
\label{eq:teacher-operator}
W^T_{ij}=\mathbf 1[j\in\mathcal M_k^T(i)]
\exp(c^T_{ij}/\tau_i),
\qquad
P^T=(D^T)^{-1}W^T,
\end{equation}
where $D^T$ is the row-degree matrix. Unlike a batch-local support, each row of $P^T$ is anchored in the teacher's corpus geometry. We use an entropy-calibrated bandwidth $\tau_i$ to accommodate variation in neighborhood density; its construction is given in Appendix~\ref{app:entropic-bandwidth}. Once $P^T$ is stored, the teacher is not used during training.

\subsection{From Neighborhoods to Multi-Scale Topology}
\label{sec:kernels}

The corpus graph corrects the supervision support, but its individual rows describe only immediate neighbors. Global geometry also depends on how these neighborhoods connect. We expose this structure with the lazy transition operator $\widetilde P^T=\tfrac12(I+P^T)$. For a selected set of scales $\mathcal R$ containing $1$, the teacher diffusion profiles are
\begin{equation}
\label{eq:diffusion-targets}
g^T_{i,1}=P_i^T,
\qquad
g^T_{i,r}
=\left.\delta_i^\top(\widetilde P^T)^r\right|_{\mathcal D\setminus\{i\}},
\quad r\in\mathcal R\setminus\{1\}.
\end{equation}
Every profile is viewed as a distribution on the common corpus index set, with zero mass at excluded entries. The one-step profile preserves direct teacher neighbors, whereas broader profiles summarize the corpus regions connected to an anchor~\citep{coifman2006diffusionmaps}. We use the collection $\{g^T_{i,r}\}_{i,r}$ as a multi-scale representation of teacher geometry.

At each scale, the ideal full-corpus student profile is the direct readout
\begin{equation}
\label{eq:student-global-profile}
g^S_{i,r}(j)
=\frac{\exp(c^S_{ij}/\tau_{i,r})}
{\sum_{\ell\neq i}\exp(c^S_{i\ell}/\tau_{i,r})},
\qquad j\neq i.
\end{equation}
This is the ideal full-corpus object approximated by the RIPPLE loss; it does not posit a student Markov walk. We measure the geometry induced by these profiles through
\begin{equation}
\label{eq:diffusion-distance}
d_r^M(i,j)=\mathrm{TV}(g^M_{i,r},g^M_{j,r}),
\qquad M\in\{T,S\}.
\end{equation}

\begin{theorem}[Diffusion-geometry preservation]
\label{thm:profile-geometry}
Fix a scale $r$. If
$\max_i D_{\mathrm{KL}}(g^T_{i,r}\|g^S_{i,r})\leq\varepsilon_r$, then
\begin{equation}
\label{eq:profile-uniform-bound}
\sup_{i,j}|d_r^T(i,j)-d_r^S(i,j)|
\leq\sqrt{2\varepsilon_r}.
\end{equation}
Moreover, for
$\bar\varepsilon_r=N^{-1}\sum_iD_{\mathrm{KL}}(g^T_{i,r}\|g^S_{i,r})$,
\begin{equation}
\label{eq:profile-average-bound}
\frac1{N^2}\sum_{i,j}|d_r^T(i,j)-d_r^S(i,j)|
\leq\sqrt{2\bar\varepsilon_r}.
\end{equation}
\end{theorem}

Theorem~\ref{thm:profile-geometry} connects the training target to the central goal: aligning anchor profiles controls distortion of the pairwise diffusion geometry. RIPPLE therefore matches teacher diffusion profiles directly at multiple resolutions. The proof is in Appendix~\ref{app:proof-profile}; the scale construction and temperatures are deferred to Appendix~\ref{app:hyper-ladder}.

\subsection{Mass-Aware Candidate Sampling}
\label{sec:candidates}

Diffusion supplies corpus-level supervision, but it also expands the support of each target. Scoring all $N$ corpus columns per anchor would eliminate the efficiency of mini-batch training. RIPPLE instead samples candidates from the teacher target mass.

\begin{lemma}[Mass coverage of target-proportional sampling]
\label{lem:sampling}
Let $p$ be a distribution and let $C_m$ contain the distinct atoms observed in $m$ independent draws from $p$. Then
\begin{equation}
\label{eq:sampling-mass}
\mathbb E[p(C_m^c)]
=\sum_jp_j(1-p_j)^m
\leq e^{-mt}+\operatorname{Tail}_p(t),
\qquad
\operatorname{Tail}_p(t)=\sum_{j:p_j<t}p_j,
\end{equation}
for every $t>0$.
\end{lemma}

Lemma~\ref{lem:sampling} shows why this approximation can remain compact: candidate coverage is governed by the concentration of the target, rather than directly by corpus size. RIPPLE applies the result at each scale, retains the highest-mass entries, and samples the remaining mass-aware quota without replacement. Appendix~\ref{app:proof-sampling} gives the proof and implementation connection.

For anchor $i$, the scale-wise draws are merged with hard and random negatives into one shared candidate set $C_i$. For each scale, define
\begin{equation}
\kappa_{i,r}
=D_{\mathrm{KL}}(g^T_{i,r}|_{C_i}\|g^S_{i,r}|_{C_i}),
\quad
\delta^T_{i,r}=g^T_{i,r}(C_i^c),
\quad
\delta^S_{i,r}=g^S_{i,r}(C_i^c).
\end{equation}

\begin{corollary}[Sampled control of diffusion geometry]
\label{cor:sampled-topology}
Let
$\eta_{i,r}=\sqrt{\kappa_{i,r}/2}+\delta^T_{i,r}+\delta^S_{i,r}$.
For any pair of anchors $i,j$,
\begin{equation}
\label{eq:sampled-topology-bound}
|d_r^T(i,j)-d_r^S(i,j)|
\leq\eta_{i,r}+\eta_{j,r}.
\end{equation}
\end{corollary}

Corollary~\ref{cor:sampled-topology} connects the sampled objective back to global geometry. Its bound separates restricted matching error, missed teacher mass, and student leakage. Target-proportional sampling reduces $\delta^T$; hard and random negatives, together with an ambient comparison term, are used to discourage $\delta^S$. We measure the latter rather than claiming that finite negative samples guarantee full-corpus leakage control.

For an anchor batch $B$, we pool and deduplicate its candidates,
\begin{equation}
\label{eq:pooled-candidates}
\mathcal P_B=\bigcup_{i\in B}C_i.
\end{equation}
Each text in $\mathcal P_B$ is encoded once, while every anchor retains its own target on its sampled support. To expose mass outside the graph-selected candidates, we also use an ambient profile
$g^M_{i,0}=\operatorname{softmax}_{j\neq i}(c^M_{ij}/\tau_0)$ for $M\in\{T,S\}$.
Let $\mathcal R_0=\{0\}\cup\mathcal R$, with $\Omega_{i,0}=\mathcal P_B\setminus\{i\}$ and $\Omega_{i,r}=C_i$ for $r\geq1$. The resulting multi-scale geometry loss is
\begin{equation}
\label{eq:ldiff}
\mathcal L_{\mathrm{diff}}
=\frac1{|B|}\sum_{i\in B}\sum_{r\in\mathcal R_0}
\omega_rD_{\mathrm{KL}}
\left(g^T_{i,r}|_{\Omega_{i,r}}\|g^S_{i,r}|_{\Omega_{i,r}}\right).
\end{equation}

\subsection{Row Supervision for Cross-Anchor Consistency}
\label{sec:walk}

Mass-aware sampling makes global-profile matching computationally feasible, but a second gap remains. A student embedding induces one symmetric cosine geometry, whereas the softmax loss normalizes every anchor row independently. Perfect row matching can therefore preserve within-row probabilities without fully coupling the underlying similarities across anchors.

\begin{proposition}[Softmax gauge and reciprocal-overlap consistency]
\label{prop:row-consistency}
Suppose teacher and student one-step rows use the same temperature on a column set $\Omega_i$ with $|\Omega_i|\geq2$. If their row softmaxes coincide, then there is a row-dependent constant $a_i$ such that
\begin{equation}
\label{eq:row-gauge}
c^S_{ij}=c^T_{ij}+a_i,
\qquad j\in\Omega_i.
\end{equation}
Construct the reciprocal-overlap graph on supervised rows, with an edge $\{u,v\}$ whenever $v\in\Omega_u$ and $u\in\Omega_v$. The offsets are constant on every connected component of this graph.
\end{proposition}

Proposition~\ref{prop:row-consistency} identifies the missing constraint: reciprocal overlap collapses independent row gauges to one component-wise constant. RIPPLE creates these constraints by following the teacher transition graph and supervising the rows encountered along its paths. Starting at an anchor, a walk visits
\begin{equation}
\label{eq:walk}
j_0=i,
\qquad
j_{t+1}\sim P^T_{j_t}.
\end{equation}
For every visited row $u$ with $|\Omega_u^w|\geq2$, let $\Omega_u^w=\mathcal P_B\cap\mathcal M_k^T(u)$ and let $\nu_B(u)$ be its normalized visit count. We add
\begin{equation}
\label{eq:lrow}
\mathcal L_{\mathrm{row}}
=\sum_u\nu_B(u)
D_{\mathrm{KL}}\!\left(
P_u^T|_{\Omega_u^w}\ \middle\|\
\operatorname{softmax}_{v\in\Omega_u^w}(c^S_{uv}/\tau_u)
\right).
\end{equation}
A realized walk edge supplies a reciprocal constraint whenever both endpoints and cross-columns are present in the supervised pool. Thus the teacher graph not only selects informative columns, but also selects the additional rows needed to connect their normalization gauges. These rows largely reuse existing encodings. Sampling details are given in Appendix~\ref{app:hyper-nonbacktracking}, and Proposition~\ref{prop:row-consistency} is proved in Appendix~\ref{app:proof-gauge}.

The two components give the RIPPLE objective
\begin{equation}
\label{eq:full-objective}
\mathcal L_{\mathrm{RIPPLE}}
=\mathcal L_{\mathrm{diff}}+\lambda\mathcal L_{\mathrm{row}}.
\end{equation}
The graph, scales, sampling budgets, temperatures, and loss weights are specified in Appendix~\ref{app:hyper}. Exact graph construction costs $O(N^2d_T)$ before sparsification and stores $O(Nk)$ edges; approximate nearest-neighbor search can replace it at larger scale. Each update encodes at most $|B|+|\mathcal P_B|$ texts. The teacher graph and all distillation machinery disappear at inference.
